\introduction % Структурный элемент: ВВЕДЕНИЕ

Распределённые системы лежат в основе современных вычислительных платформ,
облачных сервисов и транзакционных приложений. Их корректность критически
важна, поскольку ошибки в алгоритмах координации, согласования и обмена
сообщениями могут приводить к потере данных, нарушению согласованности
состояния и отказам сервисов. При этом поведение распределённых систем
существенно усложняется из-за асинхронности, недетерминированности доставки
сообщений и возможных сбоев отдельных узлов.

Формальные методы позволяют повысить уверенность в корректности распределённых
алгоритмов за счёт построения строгих математических моделей. Одним из широко
применяемых инструментов является язык спецификаций TLA+
\cite{lamport2002specifying}, предназначенный для описания систем в терминах
состояний и переходов между ними. С помощью средства проверки моделей TLC
возможно исследование пространства состояний спецификации и доказательство
инвариантов.

Однако формальная верификация спецификации не гарантирует корректность
конкретной программной реализации алгоритма. На практике возникает разрыв между
абстрактной моделью и исполняемым кодом. В данной работе рассматривается подход
к сокращению этого разрыва путём сопоставления трасс исполнения реализации со
спецификацией TLA+.

\subsection*{Актуальность темы исследования}

Несмотря на широкое распространение формальных спецификаций в проектировании
распределённых алгоритмов, проблема согласованности спецификации и реализации
остаётся актуальной. Традиционные методы тестирования не обеспечивают полного
покрытия возможных сценариев взаимодействия компонентов, особенно в условиях
недетерминированного поведения сети.

Модельная проверка позволяет анализировать абстрактные модели, однако переход
от модели к реальному коду часто сопровождается упрощениями, оптимизациями или
изменением гранулярности атомарных действий. В результате возможны расхождения
между поведением реализации и предполагаемой спецификацией.

Перспективным направлением является модельная валидация исполнения (Model-Based
Trace Checking) \cite{howard2011tracechecking}, при которой реальные трассы
работы распределённой программы проверяются на соответствие формальной
спецификации с использованием средства модельной проверки. Такой подход
позволяет обнаруживать ошибки реализации, избыточные ограничения спецификации и
несоответствие уровня атомарности.

Разработка и практическая апробация метода валидации трасс исполнения по
спецификации TLA+ представляется актуальной задачей, направленной на повышение
надёжности распределённых систем. Этот метод активно применяется в индустрии
\cite{newcombe2015aws}.

\subsection*{Цель и задачи работы}

Целью данной научно-исследовательской работы является разработка метода
модельной валидации реализации распределённого алгоритма по формальной
спецификации TLA+ на основе анализа трасс исполнения.

Для достижения поставленной цели были сформулированы следующие задачи:

\begin{enumerate}
    \item Изучить теоретические основы языка спецификаций TLA+ и средства
модельной проверки TLC.
    \item Разработать механизм инструментирования реализации для регистрации
изменений переменных спецификации.
    \item Реализовать программный прототип распределённого алгоритма двухфазной
транзакции.
    \item Реализовать процедуру сопоставления трасс программы и спецификации
с использованием TLC.
    \item Провести экспериментальную валидацию и проанализировать полученные
результаты.
\end{enumerate}

Результатом работы является прототипный фреймворк, обеспечивающий сопоставление
трасс распределённой реализации с формальной спецификацией TLA+, а также
демонстрационная реализация алгоритма двухфазной транзакции, использованная для
апробации предложенного метода.
